\documentclass[11pt]{article}
\usepackage[margin=1in]{geometry}
\usepackage{amsmath,amssymb,amsthm}
\usepackage{hyperref}

\newtheorem{theorem}{Theorem}
\newtheorem{lemma}[theorem]{Lemma}
\newtheorem{proposition}[theorem]{Proposition}
\theoremstyle{remark}
\newtheorem{remark}[theorem]{Remark}

\newcommand{\qbinom}[2]{\genfrac{[}{]}{0pt}{}{#1}{#2}_q}

\title{Disproof of Conjecture 00000007672:\\
$\gcd(F_6,F_3)=1$ is not a multiple of $F_{\gcd(6,3)}=1+q$}
\author{jilint777}
\date{2026-10-04}

\begin{document}
\sloppy
\maketitle

\begin{abstract}
Conjecture 00000007672 asserts that for the $q$-Fibonacci polynomials
$F_n(q)=\sum_j\qbinom{n-1-j}{j}q^{j^2}$, the greatest common divisor
$\gcd(F_n,F_m)$ in $\mathbb Z[q]$ equals $F_{\gcd(n,m)}$ times a correction
factor. We show that $\gcd(F_6,F_3)=1$ in $\mathbb Z[q]$, while
$F_{\gcd(6,3)}=F_3=1+q$. Since $1$ is not $(1+q)$ times any polynomial, the
claim fails. More generally the claim fails for $(3,3k)$ for every
$k\ge2$. The main counterexample is fully formalized in Lean~4 (core library
only), including the proof that $1$ is a gcd of $F_6$ and $F_3$ in
$\mathbb Z[q]$.
\end{abstract}

\section{The conjecture}

The statement (English version) reads:
\begin{quote}
\emph{Definition:} The $q$-Fibonacci polynomials are
$F_n(q)=\sum\,[n-1-j \text{ choose } j]_q\, q^{j^2}$, with Gaussian binomial
coefficients. \emph{Conjecture:} $\gcd(F_n(q),F_m(q))$ in $\mathbb Z[q]$ equals
$F_{\gcd(n,m)}(q)$ times a correction factor from a finite exceptional set $E$
depending only on primes $p\equiv\pm2 \pmod 5$, and the count of
$E\cap[1,x]$ is $O(\log\log x)$.
\end{quote}
The Chinese version says the same thing. We refute the first clause:
\begin{equation}\label{eq:clause}
\gcd(F_n,F_m)=F_{\gcd(n,m)}\cdot c_{n,m}\qquad\text{for some polynomial } c_{n,m}.
\end{equation}
The ``correction factor'' multiplies a polynomial and is supposed to yield the
gcd, which lies in $\mathbb Z[q]$. So it is a polynomial, in $\mathbb Z[q]$ or at
least in $\mathbb Q[q]$. Our argument works for $c_{n,m}\in\mathbb Q[q]$ too
(Remark~\ref{rem:Q}). If $c_{n,m}$ were allowed to be an arbitrary rational
function, \eqref{eq:clause} would hold for any polynomial in place of
$\gcd(F_n,F_m)$ and would say nothing.

\section{Definitions and small values}

The Gaussian binomial coefficient is
$\qbinom{n}{k}=\prod_{i=1}^{k}\frac{1-q^{n-k+i}}{1-q^{i}}$ for $0\le k\le n$, and
$0$ otherwise. It satisfies the $q$-Pascal rule
$\qbinom{n+1}{k+1}=\qbinom{n}{k}+q^{k+1}\qbinom{n}{k+1}$. Only the terms with
$0\le j\le n-1-j$ contribute to $F_n$. We get
\begin{align*}
F_3&=\qbinom{2}{0}+\qbinom{1}{1}q=1+q,\\
F_6&=\qbinom{5}{0}+\qbinom{4}{1}q+\qbinom{3}{2}q^4
 =1+(q+q^2+q^3+q^4)+(q^4+q^5+q^6)\\
&=1+q+q^2+q^3+2q^4+q^5+q^6 .
\end{align*}
As a sanity check, $F_n(1)=\sum_j\binom{n-1-j}{j}$ is the Fibonacci number
$1,1,2,3,5,8,\dots$.

\begin{lemma}\label{lem:values}
$F_3(-1)=0$ and $F_6(-1)=2$. Moreover $F_6=F_3\cdot(-1+2q-q^2+2q^3+q^5)+2$.
\end{lemma}
\begin{proof}
Direct computation. $F_6(-1)=1-1+1-1+2-1+1=2$.
\end{proof}

\section{\texorpdfstring{The gcd of $F_6$ and $F_3$ in $\mathbb Z[q]$}{The gcd of F6 and F3 in Z[q]}}

\begin{proposition}\label{prop:gcd}
$1$ is a greatest common divisor of $F_6$ and $F_3$ in $\mathbb Z[q]$. That is,
every $D\in\mathbb Z[q]$ dividing both $F_6$ and $F_3$ divides $1$, i.e.\
$D=\pm1$.
\end{proposition}
\begin{proof}
Write $DA=F_3=1+q$ and $DB=F_6$ with $A,B\in\mathbb Z[q]$. Degrees add in
$\mathbb Z[q]$, because the top coefficient of a product is the product of the
top coefficients and $\mathbb Z$ has no zero divisors. Hence
$\deg D+\deg A=1$, so $D$ or $A$ is constant.

If $A=a$ is constant, then $D(0)a=1$ gives $a\ne0$, and $D(-1)a=F_3(-1)=0$
gives $D(-1)=0$. But then $F_6(-1)=D(-1)B(-1)=0$, contradicting
$F_6(-1)=2$ (Lemma~\ref{lem:values}).

If $D=d$ is constant, then $d\cdot A(0)=1$ forces $d=\pm1$, so $D\cdot d=1$
and $D\mid 1$.
\end{proof}

\begin{theorem}\label{thm:main}
For $(n,m)=(6,3)$ the clause \eqref{eq:clause} fails. Conjecture 00000007672 is
false.
\end{theorem}
\begin{proof}
$\gcd(6,3)=3$ and $F_3=1+q$. By Proposition~\ref{prop:gcd} the gcd of $F_6$ and
$F_3$ in $\mathbb Z[q]$ is $\pm1$, which is unique up to units. If
$\pm1=(1+q)\,c$ for a polynomial $c$, then evaluating at $q=-1$ gives
$\pm1=0$, a contradiction. More generally, \emph{no} divisor $G$ of $F_6$ is a
multiple of $F_3$. If $G=F_3c$ and $GH=F_6$, then
$2=F_6(-1)=F_3(-1)c(-1)H(-1)=0$.
\end{proof}

\begin{remark}\label{rem:Q}
The same evaluation argument works if the correction factor $c$ is allowed to
lie in $\mathbb Q[q]$, or even to be any rational function without a pole at
$q=-1$. Over $\mathbb Q[q]$ the gcd is also $1$, because $F_6\equiv 2 \pmod{1+q}$.
\end{remark}

\begin{remark}[Infinitely many failures]
The polynomials $F_n$ satisfy Schur's recurrence
$F_n=F_{n-1}+q^{n-2}F_{n-2}$ (checked for $n\le30$ in \texttt{verify.py}). Put
$f_n=F_n(-1)$. Then $f_{2k}=f_{2k-1}+f_{2k-2}$ and $f_{2k+1}=f_{2k}-f_{2k-1}$,
so $f_{2k+1}=f_{2k-2}$ and $f_{2k}=f_{2k-2}+f_{2k-4}$. The even-indexed values
are therefore Fibonacci numbers, and $f_n>0$ for every $n\ne3$. Hence
$1+q=F_3$ divides no $F_{3k}$ with $k\ge2$, and the argument of
Theorem~\ref{thm:main} refutes \eqref{eq:clause} for every pair $(3,3k)$,
$k\ge2$. The Lean project checks $f_{3k}\ne0$ for $k=2,\dots,6$. So the failure
cannot be confined to a finite exceptional set of pairs, under any reading of
the set $E$.
\end{remark}

\section{Lean formalization}

The directory \texttt{lean4/} is a self-contained Lean~4.19.0 project that uses
only the core library (no Mathlib).
\begin{itemize}
\item Polynomials in $\mathbb Z[q]$ are coefficient lists. \texttt{PEq} means
equal coefficients, \texttt{Dvd a b} is $\exists c,\ a\cdot c=b$ in $\mathbb Z[q]$,
and \texttt{IsGcd G a b} is the standard definition of a greatest common
divisor (a common divisor that every common divisor divides).
\item \texttt{gauss} defines Gaussian binomials by the $q$-Pascal rule.
\texttt{gauss\_product\_formula} checks
$\qbinom{n}{k}(q;q)_k(q;q)_{n-k}=(q;q)_n$ for all $k\le n\le7$.
\texttt{qfib n} is $\sum_{j<n}\qbinom{n-1-j}{j}q^{j^2}$ exactly as in the
conjecture. \texttt{qfib\_at\_one} checks $F_n(1)=\mathrm{Fib}_n$ for $n\le15$.
\item \texttt{common\_divisor\_dvd\_one} and \texttt{isGcd\_one} formalize
Proposition~\ref{prop:gcd}. They use the degree argument
(\texttt{coeff\_mulL\_top}, \texttt{const\_or\_const}) and evaluation
(\texttt{evalL\_mulL}, \texttt{evalL\_congr}).
\item \texttt{GcdClause} is the clause \eqref{eq:clause}: for all $n,m\ge1$ and
every gcd $G$ of $F_n,F_m$ in $\mathbb Z[q]$, $G=F_{\gcd(n,m)}\cdot c$ for some
$c\in\mathbb Z[q]$. \texttt{conjecture\_00000007672\_false} proves
$\neg$\texttt{GcdClause}. \texttt{no\_divisor\_of\_F6\_is\_multiple\_of\_F3}
proves the stronger statement from the proof of Theorem~\ref{thm:main}.
\end{itemize}
\texttt{lake build} succeeds. There is no \texttt{sorry}, no
\texttt{native\_decide}, and no added axioms. \texttt{\#print axioms} shows at
most \texttt{propext}, \texttt{Classical.choice} and \texttt{Quot.sound}. The
numerical facts are proved by \texttt{decide}; \texttt{qfib\_mult3\_at\_neg\_one}
uses \texttt{decide +kernel}, which is checked by the Lean kernel.

\section{Reproduction}
\begin{verbatim}
cd lean4 && lake build       # Lean 4.19.0, core only
python3 verify.py            # standard library only
pdflatex report.tex
\end{verbatim}
\texttt{verify.py} computes Gaussian binomials from the product formula by exact
division, which is independent of the $q$-Pascal rule used in Lean. It
recomputes $F_n$, finds $\gcd(F_6,F_3)=1$ by the Euclidean algorithm in
$\mathbb Q[q]$ (and in $\mathbb Z[q]$ by Gauss's lemma, since $F_3$ is
primitive), and lists all $76$ pairs $n<m\le30$ for which
$F_{\gcd(n,m)}\nmid F_n$ or $F_{\gcd(n,m)}\nmid F_m$.

\begin{thebibliography}{9}
\bibitem{carlitz} L.~Carlitz, Fibonacci notes 3: $q$-Fibonacci numbers,
\emph{Fibonacci Quart.} 12 (1974), 317--322.
\bibitem{andrews} G.~E.~Andrews, Fibonacci numbers and the Rogers--Ramanujan
identities, \emph{Fibonacci Quart.} 42 (2004), 3--19.
\end{thebibliography}

\end{document}
